\chapter*{Preface: What These Notes Are For}
\addcontentsline{toc}{chapter}{Preface: What These Notes Are For}

\begin{remarkbox}
These notes are designed as a continuation of the Classical Field Theory I--II
sequence. They do not try to re-teach variational calculus, Lorentz covariance,
Noether currents, gauge redundancy, or the classical equations of fundamental
fields from the beginning. Instead, they start where the classical notes stop:
with fields as dynamical systems ready to be quantized.
\end{remarkbox}

\section*{Quantum Field Theory After Classical Field Theory}

Classical field theory teaches us to describe nature by local fields, actions,
symmetries, conservation laws, and equations of motion. Quantum field theory
keeps this architecture, but changes the meaning of the basic objects. A
classical field configuration is no longer the final description of a physical
state. It becomes an ingredient in a quantum theory whose observable content is
encoded in states, operators, correlation functions, transition amplitudes, and
probabilities.

The guiding idea of the first part of these notes is simple: a free classical
field decomposes into normal modes, and each normal mode behaves like a harmonic
oscillator. Quantizing the field therefore means quantizing infinitely many
oscillators at once. The excitations of these oscillators are interpreted as
particles.

\section*{What Will Not Be Repeated}

The notes assume familiarity with the following ideas from classical field
theory:
\begin{itemize}
    \item the action principle and Euler--Lagrange equations for fields,
    \item canonical momenta and Hamiltonian density,
    \item Lorentz covariance and relativistic notation,
    \item real and complex scalar fields,
    \item electromagnetic and Yang--Mills gauge fields at the classical level,
    \item Noether's theorem and conserved currents,
    \item Green's functions as solutions of linear differential equations.
\end{itemize}

Whenever one of these topics is used in a quantum calculation, the notes recall
just enough of the classical structure to make the calculation readable. The
aim is not repetition, but conversion: classical structures are translated into
quantum structures.

\section*{What Is New in QFT}

The genuinely new concepts are the vacuum, Fock space, field operators,
propagators, Wick contractions, Feynman diagrams, scattering amplitudes, loop
corrections, regularization, renormalization, and scale dependence. These ideas
are not decorative additions to classical field theory. They are the mechanisms
that make quantum field theory predictive.

A central conceptual warning will appear repeatedly: a Feynman diagram is not a
literal picture of tiny particles following tiny trajectories. It is a compact
bookkeeping device for terms in a perturbative expansion of correlation
functions or scattering amplitudes.

\section*{How to Read These Notes}

A useful rhythm for reading a QFT calculation is:
\[
    \text{classical field}
    \longrightarrow
    \text{mode expansion}
    \longrightarrow
    \text{operators}
    \longrightarrow
    \text{states}
    \longrightarrow
    \text{correlators}
    \longrightarrow
    \text{amplitudes}.
\]

At first this may look like a change of notation. It is not. Each arrow changes
what is meant by a physical prediction. The purpose of the notes is to make
those changes explicit and technically usable.
